\documentclass[11pt]{scrartcl}

\usepackage[stixtwo]{fontsetup}
\usepackage{amsmath}

\begin{document}
\[
	c_n = \frac{-6 \cdot (-6)^n + \frac{n^3 4^n}{2}}{\frac{3}{2} \cdot 6^n + n}
\]
\[
	c_{2n}
	= \frac{-6 \cdot 36^n + 4n^3 16^n}{1{,}5 \cdot 36^n + 2n}
	= \frac{-6 + 4n^3 \left(\frac{4}{9}\right)^n}{1{,}5 + \frac{2n}{36^n}}
	\to -4
\]
\[
	c_{2n - 1}
	= \frac{36^n + \frac{(2n - 1)^3 16^n}{8}}{0{,}25 \cdot 36^n + 2n - 1}
	= \frac{1 + \frac{(2n - 1)^3}{8} \left(\frac{4}{9}\right)^n}{0{,}25 + \frac{2n - 1}{36^n}}
	\to 4
\]
Minden elem szerepel valamelyik részsorozatban, így 2 torlódási pont van:
$\pm 4$. $\limsup c_n = 4$, $\liminf c_n = -4$, a határérték nem létezik.

\end{document}
